\begin{lemma}[Scheide maximal-tree package (a)]
Let $(G,e_0,\varphi)$ be a $k$-triple and let $T$ be a fixed-pair maximal
Tashkinov tree for $(e_0,\varphi)$.  Then $V(T)$ is elementary and closed
with respect to $\varphi$.
\end{lemma}
\begin{proof}
The Tashkinov-tree clause in
\noderef{FixedPairMaximalTashkinovTree}, together with the $k$-triple
hypothesis, lets us apply \noderef{TashTree}; hence $V(T)$ is elementary.

Suppose $V(T)$ were not closed.  By \noderef{ColoredVertexSetClosed} and
\noderef{BoundaryEdges}, there would be a boundary edge identity $f$ with
one end $x\in V(T)$ and the other end $y\notin V(T)$, and a color $\eta$
such that $\varphi(f)=\eta$ and $\eta$ is missing at some vertex
$z\in V(T)$.  Append $f,y$ to the displayed tree-sequence $T$.  The new
vertex is distinct from all old vertices.  The boundary identity $f$ is
distinct from every old tree edge, since each old tree edge has both ends
in $V(T)$ by \noderef{TreeSequence}, whereas $f$ has its other end $y$
outside.  Thus the appended sequence is a tree-sequence.  Its new colored
edge has the earlier missing-color witness $z$, so it is a Tashkinov tree
by \noderef{TashkinovTree}.  Its vertex set strictly contains $V(T)$,
contradicting fixed-pair maximality.  Therefore $V(T)$ is closed as well.
\end{proof}
